\begin{figure}[H]
\centering
\TeachingFlow{Initial cost K}{Savings M minus A per change}{Break-even after K/(M−A) changes}
\caption{The planning relation requires positive savings on comparable changes.}\label{fig:teach-P-04-prototype-06}
\end{figure}

For synthetic costs K = 100,000, M = 2,000 and A = 1,500 in the same currency, the break-even count is 200 comparable changes. If assisted cost rises to 2,100, this labor-saving calculation has no positive break-even point. Avoided errors would require separate evidence and valuation.

